% Part: lambda-calculus
% Chapter: syntax
% Section: abbreviated-syntax

\documentclass[../../../include/open-logic-section]{subfiles}

\begin{document}

\olfileid{lam}{syn}{abb}
\olsection{مختصر نحو}

\olref[trm]{defn:term} وچ تعریف کیتے ٹرم کدے لکھنے اوکھے
ہوندے نیں، ایس لئی ہور مختصر نحو ورتنی سوکھی اے۔ پر ایہ دھیان
رکھنا اے کہ مختصر لکھت وچ وی ہر ٹرم اکّو طرح پڑھیا جا سکے۔
وڈی پدھر اُتے ورتی جاندی اک صورت \emph{مختصر ٹرم} اے، جیہڑی
ایویں لکھی جاندی اے۔

\begin{enumerate}
\item جدوں قوساں چھڈ دئیے، اطلاق کھبے توں سجے ول جُڑدا اے۔
  مثلاً جے $M$، $N$، $P$ تے~$Q$ ٹرم نیں، تاں $MNPQ$،
  $(((MN)P)Q)$ دا مختصر روپ اے۔
\item قوساں چھڈن اُتے لامبڈا تجرید دا دائرۂ اثر جِنّا ہو سکے
  وڈا رکھدے آں۔ مثلاً $\lambd[x][MNP]$ نوں
  $(\lambd[x][MNP])$ پڑھدے آں۔
\item اک لامبڈا نال کئی متغیراں دی تجرید لکھ سکدے آں۔
  مثلاً $\lambd[xyz][M]$،
  $\lambd[x][\lambd[y][\lambd[z][M]]]$ دا مختصر روپ اے۔
\end{enumerate}

مثلاً،
\[
\lambd[xy][xxyx \lambd[z][xz]]
\]
ایہ ہِٹھلی عبارت دا مختصر روپ اے:
\[
(\lambd[x][(\lambd[y][((((xx)y)x)(\lambd[z][(xz)]))])]).
\]

\begin{prob}
مختصر ٹرم $\lambd[g][(\lambd[x][g (x x)])
  \lambd[x][g (x x)]]$ نوں پورا لکھو۔
\end{prob}

\end{document}
